\section{经济与社会：霸权背后的土地、劳役与交通}

春秋战争并不是只靠勇士决斗。国家要长期会盟和出兵，需要粮食、车马、道路、城邑和劳役。大国能够维持更大规模的军事与外交，往往因为腹地、河流、平原和人口动员更有优势；小国则更依赖城防、外交和在大国之间选择时机。

一场会盟看上去发生在诸侯帐前，真正的成本却落在更远的地方。车从哪里来，粮怎样送到集结地，谁离开田地服役，城邑如何在主力出征时继续守住？传世史书很少逐笔记录这些账目，但会盟、筑城、救援和长期战争反复出现，说明政治秩序始终由大量看不见的运输、供给和劳役托着。

春秋中后期，贵族内部的家族与封邑关系也在改变。卿大夫掌握的资源越来越多，国君与卿族之间的冲突逐渐成为各国政治的核心。晋的卿族、齐的田氏、鲁的三桓，表面上是不同国家的家族故事，背后却是同一个问题：原本替国君经营地方和军队的人，如何逐渐把这些资源变成了自己的政治基础。

对小国而言，生存并不只是选边站。郑靠近周王室，能在特定时刻挑战王师；宋能在晋楚之间提供会盟场所；鲁留下年表与礼仪传统，却也被国内家族政治牵制。小国的脆弱不意味着没有行动能力，只是它们的行动通常必须借大国均势和地理位置才有空间。

\begin{turningpoint}[从春秋社会到战国国家]
春秋的资源仍多通过宗族、封邑和礼仪关系流动。战国改革要做的，正是把土地、户籍、军役和税收更直接地接到国家手中。这会提高动员效率，也会改变普通家庭与国家发生关系的方式。春秋的卿族问题、会盟成本和小国困境，都是战国制度竞争的社会前史。
\end{turningpoint}
